\documentclass[11pt]{article}
\usepackage[a4paper,margin=22mm]{geometry}
\usepackage{amsmath,amssymb,mathtools,amsthm}
\usepackage{booktabs}
\usepackage{enumitem}
\usepackage{xcolor}
\usepackage[hidelinks]{hyperref}

\newcommand{\R}{\mathbb{R}}
\newcommand{\one}{\mathbf 1}
\newcommand{\Var}{\operatorname{Var}}
\newcommand{\Cov}{\operatorname{Cov}}
\newcommand{\sg}{\operatorname{stopgrad}}
\newcommand{\rnd}{\operatorname{round}}
\newtheorem*{claim}{Claim}

\setlist{itemsep=2pt,topsep=3pt}
\setlength{\parindent}{0pt}
\setlength{\parskip}{5pt}

\newcounter{qn}
\newcommand{\Q}[1]{\refstepcounter{qn}\par\medskip\noindent\textbf{Q\theqn.\ #1}\par\nopagebreak}
\newcommand{\A}{\noindent\textit{Answer.}\ }

\begin{document}

\begin{center}
{\LARGE\bfseries Revision notes}\par\vspace{4pt}
{\large Ye et al., ``Robust Training of Neural Networks at Arbitrary Precision and Sparsity'' (Google DeepMind, ICLR 2026)}\par\vspace{4pt}
Ayush Srivastava (2026AIB2015) \quad Sarvesh Nikas (2026AIB2045)
\end{center}

\tableofcontents

\section{Notation}

A \emph{block} is a vector $x\in\R^N$ that shares one scale and one shift. Means and variances are population moments ($1/N$):
\[
\bar x=\tfrac1N\one^\top x,\qquad
\Var_q=\tfrac1N\|q-\bar q\one\|^2,\qquad
\Cov_{xq}=\tfrac1N(x-\bar x\one)^\top(q-\bar q\one).
\]
Throughout, $u=q-\bar q\one$ (centred integers), $k$ is the bit width, and the integer grid is $\{0,1,\dots,2^k-1\}$.

\section{Why quantize, and what training looks like}

\paragraph{Motivation.} A language model has billions of weights stored in 16 or 32 bits. Quantization stores each in $k\in\{4,2,1\}$ bits, plus one scale and shift per block. This cuts memory, memory traffic and arithmetic cost.

\paragraph{Weights vs.\ activations.} \emph{Weights} are learned parameters, fixed after training. \emph{Activations} are the vectors flowing between layers for a given input; they change with every input. A linear layer computes $y=Wx$ with $W$ weights and $x$ activations. Both are quantized because the cost of a multiply scales with (weight bits)$\times$(activation bits): 1-bit weights with 16-bit activations still need 16-bit additions.

\paragraph{Quantization-aware training (``fake quantization'').} Training keeps the float weights $x$ and updates them by gradient descent. The forward pass uses the quantized-then-reconstructed values, so the network learns to tolerate rounding. Reconstruction is done in float during training only so that autograd can run; it carries exactly the information of the integers $q$. After training only $q$ and the per-block scalars are kept.

\paragraph{Where quantization is applied.} At the input and weights of every linear layer (in a transformer: QKV and output projections, MLP up and down projections). Embeddings, LayerNorm, softmax and the attention products $QK^\top$, $\mathrm{softmax}\cdot V$ stay in float in our runs.

\section{The quantizer (paper Eq.~1)}

\[
x\xrightarrow{\ f\ }f(x)\xrightarrow{\ \rnd\ }q\xrightarrow{\ s_fq+b_f\one\ }\hat x,
\qquad
f(x)=\frac{x-b_f\one}{s_f},\quad s_f=\frac{x_{\max}-x_{\min}}{2^k-1},\quad b_f=x_{\min}.
\]
Write $q=f(x)+\delta$ with rounding error $|\delta_i|\le\tfrac12$. Then
\[
\boxed{\hat x=s_fq+b_f\one=x+s_f\delta.}
\]
Quantization is $x$ plus a scaled rounding error. Everything else builds on this identity.

\paragraph{Worked example ($k=2$).} $x=(0.9,-0.3,0.1,0.6)$. Then $s_f=1.2/3=0.4$, $b_f=-0.3$, $f(x)=(3,0,1,2.25)$, $q=(3,0,1,2)$, $\hat x=(0.9,-0.3,0.1,0.5)$. Squared error $\|\hat x-x\|^2=0.010$.

\section{The problem: the true gradient is zero}

By the chain rule, with $s_f,b_f$ held fixed,
\[
\frac{\partial\hat x}{\partial x}=s_f\cdot\frac{\partial q}{\partial f(x)}\cdot\frac1{s_f}=\frac{\partial q}{\partial f(x)}.
\]
$\rnd$ is a step function: its derivative is $0$ off the half-integers and undefined on them. So $\nabla_xL=0$ almost everywhere. A small change of $x$ does not change $q$; $x$ and all earlier layers never move.

\textit{Caveat.} $s_f$ and $b_f$ depend on $x$ through $x_{\max}$ and $x_{\min}$, so the two extreme entries of each block do get a non-zero gradient. The paper's analysis (and ours) holds $s_f,b_f$ fixed.

\section{The standard fix: the straight-through estimator (STE)}

Detach the rounding error:
\[
q=f(x)+\delta,\qquad \delta=\sg\big(\rnd(f(x))-f(x)\big).
\]
$\sg$ is the identity in the forward pass and has zero derivative in the backward pass. Hence $\partial q/\partial f(x)=I$ and
\[
\text{forward: }\hat x=x+s_f\delta,\qquad \text{backward: }\frac{\partial\hat x}{\partial x}=I,\quad \nabla_xL=\nabla_{\hat x}L.
\]
(Bengio, L\'eonard, Courville, 2013.)

\paragraph{The blind spot.} The forward pass is shifted by $s_f\delta$, but the Jacobian is $I$ for every $\delta$: the gradient is exactly what it would be with no quantization. Earlier layers receive no signal about how rounding hurts the loss. At 8 bits $s_f\delta$ is tiny. At 1--2 bits it is as large as the values themselves.

\paragraph{Paper's evidence (Fig.~1).} (a) Small GPT on Tiny Shakespeare, A1W1: STE's loss climbs and oscillates around 5--8; the new method converges. (b) Gemma 3 1B, A1W1, SCQ128: under STE affine quantization gives nothing (linear 0.3399, affine 0.3397); with $g$ it helps (0.3547 vs.\ 0.3751). The paper's reading: STE cannot learn the sensitive shift term, because its gradient ignores the rounding error.

\section{The new method: ridge dequantization $g$ (Eqs.~2--3)}

Keep $f$ and $\rnd$. Replace $\hat x=s_fq+b_f\one$ by an affine map whose coefficients are \emph{fitted to $x$}:
\[
g(q)=s_gq+b_g\one,\qquad
(s_g,b_g)=\arg\min_{s,b}\ \frac1{2N}\|sq+b\one-x\|^2+\frac\lambda2s^2 .
\]

\begin{claim}
$s_g=\dfrac{\Cov_{xq}}{\Var_q+\lambda}$, \quad $b_g=\bar x-s_g\bar q$.
\end{claim}
\begin{proof}
The objective $J(s,b)$ is a strictly convex quadratic for $\lambda>0$. Setting $\partial J/\partial b=\frac1N\sum_i(sq_i+b-x_i)=0$ gives $b=\bar x-s\bar q$. Substituting, the residual is $s\,u_i-(x_i-\bar x)$, and
$\partial J/\partial s=\frac1N\sum_i\big(s u_i-(x_i-\bar x)\big)u_i+\lambda s=s\Var_q-\Cov_{xq}+\lambda s=0$.
(Using $\sum_iu_i=0$ to drop the $b$ terms.)
\end{proof}

Equivalent centred form: $g(q)=s_g(q-\bar q\one)+\bar x\one$.

\paragraph{Role of $\lambda$ (the ``denoising knob'', default $0.01$ everywhere).}
\begin{itemize}
\item $\lambda\to\infty$: $s_g\to0$, $g\to\bar x\one$ (ignore the noisy $q$, output the mean).
\item $\lambda=0$: ordinary least squares. Undefined when $\Var_q=0$ (all entries round to the same integer); Table~13 shows $\lambda=0$ gives NaN.
\item Table~13 (A1W1, WMT EN--DE BLEU): $\lambda=1$: 5.83; $0.01$: 21.42; $10^{-4}$: 20.08; $0$: NaN.
\end{itemize}

\paragraph{$\delta$ is inside $s_g$.} From $x=s_f(q-\delta)+b_f\one$,
\[
\Cov_{xq}=s_f\big(\Var_q-\Cov_{q\delta}\big)\quad\Longrightarrow\quad
s_g=s_f\,\frac{\Var_q-\Cov_{q\delta}}{\Var_q+\lambda}.
\]
The residual is
\[
r=x-g(q)=(s_f-s_g)\,u-s_f(\delta-\bar\delta\one),
\]
i.e.\ the rounding error left over after the fit.

\paragraph{Is $g$ a better reconstruction?}
\begin{itemize}
\item With $\lambda=0$, yes always: $\hat x$ is one particular affine function of $q$, and $g$ is the least-squares best one, so $\|g(q)-x\|\le\|\hat x-x\|$.
\item With $\lambda>0$, usually but not always (the penalty can push it slightly off). Our check on $10^4$ random blocks: $g$ is worse in 153 blocks at 1 bit; median error ratio $g/\hat x$ is 0.21 at 1 bit and 0.97 at 4 bits.
\item Worked example: $\bar q=1.5$, $\Var_q=1.25$, $\Cov_{xq}=0.5125$, $s_g=0.4067$, $b_g=-0.2851$, $g(q)=(0.935,-0.285,0.122,0.528)$, squared error $0.007$ vs.\ $0.010$.
\end{itemize}

\paragraph{The 1-bit case.} With $q\in\{0,1\}$, an affine map of $q$ just picks two levels. Least squares ($\lambda=0$) puts them at the \emph{means of the two groups}. STE's $\hat x$ puts them at $x_{\min}$ and $x_{\max}$, which is too far out. This is why STE's activations grow too large at 1 bit.

\paragraph{Linear (symmetric) variant (Eq.~4).} For zero-centred data use $g(q)=s_gq$ with $s_g=\overline{qx}/(\overline{q^2}+\lambda)$. For signs $q=\pm1$: $s_g=\overline{|x|}/(1+\lambda)$, close to the mean absolute value (the XNOR-Net scale). STE with $s_f=\max|x|$ uses levels $\pm\max|x|$.

\paragraph{Which variant where.} Affine $g$ for uncentred data such as activations (Sec.~3.3.1); linear $g$ for roughly zero-centred data, e.g.\ ternary sparse weights.

\section{$g$ as a normalization layer (Sec.~4)}

\[
g(q)=\underbrace{\frac{\Cov_{xq}}{\sqrt{\Var_q+\lambda}}}_{\sim\gamma}\cdot
\underbrace{\frac{q-\bar q\one}{\sqrt{\Var_q+\lambda}}}_{\sim(q-\mu)/\sigma}
+\underbrace{\bar x\one}_{\sim\beta}.
\]
This is LayerNorm applied to $q$, with $\lambda$ in the role of $\epsilon$, and with $\gamma,\beta$ \emph{computed} from $x$ rather than learned. Cost is about one LayerNorm per block (RMSNorm for the linear variant). $g$ has no learnable parameters.

\section{The backward pass: how $\delta$ reaches the gradient}

Both methods use the same surrogate for $\rnd$: $\partial q/\partial x=I/s_f$. Only the step after $\rnd$ differs.

\paragraph{STE.} $\partial\hat x/\partial q=s_fI$ for every $\delta$, so the Jacobian is $I$.

\paragraph{With $g$: two paths.} $g$ depends on $x$ through $q$ and also directly (through $\bar x$ and $\Cov_{xq}$):
\[
\nabla_xL=\underbrace{\frac1{s_f}\Big(\frac{\partial g}{\partial q}\Big)^{\!\top}\nabla_gL}_{\text{through }q\text{ (surrogate)}}+\underbrace{\Big(\frac{\partial g}{\partial x}\Big)^{\!\top}_{q\text{ fixed}}\nabla_gL}_{\text{through }x\text{ (exact)}}.
\]

\paragraph{Our derivation of the full Jacobian.} With $\rho=s_g/s_f$, $c=1/\big(s_fN(\Var_q+\lambda)\big)$ and
\[
P=\frac{\one\one^\top}{N}+\frac{uu^\top}{N(\Var_q+\lambda)},
\qquad
\boxed{J_g=\frac{dg}{dx}=P+\rho\,(I-P)+c\,u\,r^\top},\qquad J_{\text{STE}}=I.
\]
\begin{itemize}
\item $P$ is the exact (a.e.) Jacobian: with $q$ fixed, $g=s_gu+\bar x\one$ is linear in $x$; $\partial\bar x/\partial x=\one^\top/N$ and $\partial s_g/\partial x=u^\top/(N(\Var_q+\lambda))$. It depends only on $q$, so it carries no information about $\delta$.
\item With $\lambda=0$, $P$ is the orthogonal projection onto $\operatorname{span}\{\one,u\}$ (since $\|u\|^2=N\Var_q$ and $\one\perp u$). With $\lambda>0$ the $u$ direction is slightly shrunk.
\item $\rho(I-P)$ and $c\,ur^\top$ pass through the surrogate. $\rho$ depends on $\delta$ through $s_g$; $r$ is the leftover rounding error. This is how $\delta$ enters the gradient.
\item Sanity check: if $\rho=1$ and $r=0$, $J_g=I$, i.e.\ STE.
\item Verified against autograd in \texttt{implementation/jacobian\_check.py} (excluding the argmin and argmax entries, since $s_f,b_f$ depend on them).
\end{itemize}

\paragraph{Sizes on a random block of 64.}
\begin{center}
\begin{tabular}{lccc}
\toprule
& 1 bit & 2 bits & 4 bits\\
\midrule
$\rho=s_g/s_f$ & 0.29 & 0.87 & 0.99\\
share of upstream gradient in $c\,ur^\top$ & 0.5\% & 4.7\% & 0.3\%\\
share in exact part $P$ & 14.8\% & 21.1\% & 3.2\%\\
\bottomrule
\end{tabular}
\end{center}
At 4 bits $J_g\approx I$: $g$ behaves almost like STE, consistent with the tiny 4-bit gain. At 1 bit the surrogate path is scaled down to $0.29$ and the exact path matters.

\paragraph{The paper's own analysis (A.4)} derives only $\partial g/\partial q$, for the linear case. It misses the exact direct path $P$.

\section{Efficient affine matrix multiplication (Theorem~1, Sec.~5)}

Each row of $X$ (one token) and each column of $W$ (one output channel) has its own affine reconstruction $s(q-\bar q\one)+\bar x\one$. Then
\[
\boxed{Y=(s_Xs_W^\top)\odot\big(Q_XQ_W-n\,\bar q_X\bar q_W^\top\big)+n\,\bar x\,\bar w^\top.}
\]
\begin{proof}[Proof sketch]
Write $\hat X=D_X(Q_X-\bar q_X\one^\top)+\bar x\one^\top$ and $\hat W=(Q_W-\one\bar q_W^\top)D_W+\one\bar w^\top$. In $\hat X\hat W$ the cross terms vanish because centred rows and columns sum to zero: $(Q_X-\bar q_X\one^\top)\one=0$ and $\one^\top(Q_W-\one\bar q_W^\top)=0$. The constant term is $\bar x\one^\top\one\bar w^\top=n\bar x\bar w^\top$. The main term expands, using $Q_X\one=n\bar q_X$ and $\one^\top Q_W=n\bar q_W^\top$, to $Q_XQ_W-n\bar q_X\bar q_W^\top$.
\end{proof}
So an affine-quantized matmul is one \emph{integer} matmul plus rank-1 corrections. At 1 bit the integer matmul is XNOR plus popcount. (Paper typo: Sec.~5 text writes $\bar q_X\bar q_X^\top$.)

\paragraph{Inference cost of $g$.} For weights, $s_g,b_g$ are computed once after training. For activations, $x$ is new for every input, so each token needs its float $\bar x$ and $\Cov_{xq}$ at run time: an $O(N)$ pass, about a LayerNorm, cheap next to the $O(N\cdot\text{width})$ matmul, but not counted in the energy proxy.

\section{Sparsity as a special form of quantization (Sec.~3.4, A.1.5)}

\begin{enumerate}
\item Threshold small values to zero (e.g.\ 2:4 structured: 2 non-zeros per group of 4). Sparsity error $\delta_S=\operatorname{threshold}(x)-x$, so $x_S=x+\delta_S$.
\item Quantize: $q=f(x_S)+\delta_Q$.
\item Reconstruct with $g$, fitted against the original \emph{dense} $x$. So $s_g,b_g$ adapt to the combined error $\delta_S+\delta_Q$, and the backward pass sees both.
\end{enumerate}
1-bit weights with 2:4 sparsity give ternary weights $\{-1,0,1\}$. These are roughly centred, so the paper uses \emph{linear} $g$; this also keeps zeros at zero (affine $g$ would map 0 to $b_g\neq0$).

Costs: 2:4 halves the multiplies (energy factor 0.5). Storage: 1:4 gives $(1+2)/4=0.75$ bits per element; 2:4 needs a 4-bit mask, $(2+4)/4=1.5$.

\section{Results from the paper}

\begin{center}
\begin{tabular}{lccc}
\toprule
Gemma 3 1B, affine SCQ128 & STE & with $g$ & gain\\
\midrule
4 bits & 0.4254 & 0.4274 & $+0.002$\\
2 bits & 0.3794 & 0.4027 & $+0.023$\\
1 bit & 0.3397 & 0.3751 & $+0.035$\\
\bottomrule
\end{tabular}
\end{center}
Score: next-token accuracy. Fewer bits, bigger gain. SCQ128 = sub-channel blocks of 128.
\begin{itemize}
\item A1.5W1.5 (ternary, $\log_23\approx1.58$ bits), linear channel-wise: STE NaN, $g$ 0.3297.
\item Table~6: Gemma 3 4B, A4W1, 2:4 sparse: 650\,MB, 0.4517. Gemma 3 1B BF16: 1400\,MB, 0.4494. Gemma 3 4B BF16: 0.4706.
\item Table~7: ResNet-50 A4W4 76.45\% vs.\ float 76.41\%. Table~9: WMT EN--DE A4W4 29.71 vs.\ 29.49 BLEU.
\item Energy proxy (A.1.6): sparsity $\times$ act.\ bits $\times$ weight bits $\times$ ops. BF16 $=256$, A4W4 $=16$, A1W1 $=1$. An estimate, not a measurement. (Paper typo: ``$16\times16=26$''.)
\end{itemize}

\section{Our implementation}

\paragraph{Data.} Tiny Shakespeare: about 1.1M characters of Shakespeare play dialogue, character-level, vocabulary of 65 characters, first 90\% train, last 10\% validation. Loss is cross-entropy in nats per character; a uniform guess gives $\ln65\approx4.17$.

\paragraph{Small toy.} 4 blocks, 0.82M parameters, context 64, 3000 steps. Final validation loss, STE vs.\ $g$: 1 bit 2.123 vs.\ 2.071; 2 bits 1.812 vs.\ 1.790; 4 bits 1.592 vs.\ 1.597. Float 1.586. No run diverged.

\paragraph{Paper-size 1-bit run (\texttt{onebit/}).} nanoGPT \texttt{train\_shakespeare\_char}: 6 layers, 6 heads, width 384, 10.7M parameters. Every linear layer in the blocks has 1-bit weights (one block per output row) and 1-bit inputs (one block per token), so $N=384$ or 1536. We used batch 32, context 128, 1000--2000 steps (paper: 64, 256, 5000).
\begin{itemize}
\item Three settings: affine 1-bit, sign (linear) 1-bit, affine at $3\times$ learning rate. In all three STE stalls near 2.6 while $g$ keeps learning to about 2.35.
\item STE's residual stream is about $340\times$ too large after 1000 steps (617 vs.\ 1.8). Cause: STE's two levels sit at $x_{\min},x_{\max}$; $g$'s sit near the group means. LayerNorm hides the size from the loss, but each layer's update is swamped.
\item At lr $3\times10^{-3}$ STE's training loss turns upward after about step 75; $g$ reaches 2.29.
\item Not reproduced: the paper's loss explosion to 6--8. Likely because we trained on fewer tokens per step and fewer steps.
\end{itemize}

\section{Weak points of the paper}

\begin{enumerate}
\item ``Without gradient estimation'' (abstract, Table~8) is too strong: $\delta$ is detached, so $\rnd$ still uses the identity surrogate, as in STE. The new information about $\delta$ arrives \emph{through} that surrogate path.
\item ``For the first time'' clashes with the abstract's own ``well-understood mismatch''. LSQ (Esser et al., 2019) already puts $\delta$ into the step-size gradient (not the upstream gradient).
\item A.4 derives only $\partial g/\partial q$ for the linear case; it misses the exact path $P$.
\item Effective bits per element ignore per-block scale and shift (dense 1-bit listed as 1.00 BPE, Table~5).
\item Single runs, no error bars. The 4-bit gain ($+0.002$) may be noise.
\item BitNet is run at A1W1, though the paper itself says BitNet's lowest activation precision is 8 bits.
\item Energy savings are a proxy; no hardware or speed numbers. The per-token float cost of $g$ on activations is not counted.
\item No ablation of sparsity under STE vs.\ $g$.
\item Typos: $\bar q_X\bar q_X^\top$ in Sec.~5; $16\times16=26$ in A.1.6.
\end{enumerate}

\section{Viva questions and answers}

\Q{Why can't we just backpropagate through the quantizer?}
\A $\rnd$ is piecewise constant, so its derivative is 0 almost everywhere and undefined at half-integers. The chain rule gives $\nabla_xL=0$, so no weight before the quantizer ever moves.

\Q{What exactly is STE, and why ``straight-through''?}
\A It writes $q=f(x)+\sg(\rnd(f(x))-f(x))$. Forward it rounds; backward $\partial q/\partial f=I$, so the gradient passes through the rounding step as if it were the identity.

\Q{What is wrong with STE if training still runs?}
\A Its Jacobian is $I$ for every $\delta$, so the gradient is blind to the rounding error. Earlier layers cannot learn to be robust to rounding. Harmless at 8 bits; at 1--2 bits $s_f\delta$ is as large as the values, and STE diverges or stalls.

\Q{Does the new method avoid the surrogate gradient?}
\A No. $\delta$ is still detached, so $\partial q/\partial f=I$ exactly as in STE. What changes is the step after rounding: $g$'s Jacobian depends on $\delta$ (through $s_g$ and the residual $r$) and adds an exact direct path $P$. So ``without gradient estimation'' is an overstatement.

\Q{Derive $s_g$ and $b_g$.}
\A Set the partial derivatives of $\frac1{2N}\|sq+b\one-x\|^2+\frac\lambda2s^2$ to zero. $\partial_b$: $b=\bar x-s\bar q$. Substitute; $\partial_s$: $s(\Var_q+\lambda)=\Cov_{xq}$.

\Q{Why ridge and not plain least squares?}
\A Two reasons. (1) Stability: if all entries round to the same integer, $\Var_q=0$ and OLS divides by zero; $\lambda=0$ gives NaN in Table~13. (2) Denoising: $\lambda$ shrinks $s_g$ toward 0 when $q$ is a poor predictor of $x$, falling back to $\bar x$.

\Q{Why is only $s_g$ penalised, not $b_g$?}
\A $b_g$ is the shift; penalising it would bias the output mean away from $\bar x$. The noise comes in through $q$, which only $s_g$ multiplies.

\Q{What happens as $\lambda\to0$ and $\lambda\to\infty$?}
\A $\lambda\to0$: OLS, best affine fit, unstable when $\Var_q\approx0$. $\lambda\to\infty$: $s_g\to0$, $g(q)\to\bar x\one$, all information in $q$ discarded.

\Q{Is $g(q)$ always closer to $x$ than $\hat x$?}
\A For $\lambda=0$, yes: $\hat x$ is one affine function of $q$ and $g$ is the best one. For $\lambda>0$, not always: worse in 153 of $10^4$ random 1-bit blocks in our check, but the median error ratio is 0.21 at 1 bit.

\Q{Where does $\delta$ appear in the backward pass?}
\A In $J_g=P+\rho(I-P)+c\,ur^\top$. $\rho=s_g/s_f=(\Var_q-\Cov_{q\delta})/(\Var_q+\lambda)$ depends on $\delta$, and $r=(s_f-s_g)u-s_f(\delta-\bar\delta\one)$ is the leftover rounding error.

\Q{What is $P$?}
\A The exact Jacobian of $g$ with respect to its direct dependence on $x$ (through $\bar x$ and $\Cov_{xq}$), with $q$ fixed. For $\lambda=0$ it is the orthogonal projection onto $\operatorname{span}\{\one,u\}$. It needs no surrogate but carries no information about $\delta$.

\Q{Isn't $g$ just a rescaled STE?}
\A Partly. The surrogate path is scaled by $\rho$ rather than 1. But $g$ adds the exact path $P$ and the residual term $c\,ur^\top$, and $\rho$ itself depends on $\delta$. STE's slope never does.

\Q{Why is the 4-bit gain so small?}
\A At 4 bits $\delta$ is small relative to the spread of $q$, so $\rho\approx0.99$ and $r\approx0$; $J_g\approx I$, i.e.\ almost STE. Also single runs, so $+0.002$ may be noise.

\Q{Why does affine quantization not help under STE?}
\A The paper's explanation: the shift is a sensitive parameter, and STE's gradient ignores rounding error, so it cannot learn a useful shift. With $g$, $b_g=\bar x-s_g\bar q$ is fitted in closed form, and affine helps (0.3547 to 0.3751).

\Q{Explain the LayerNorm analogy.}
\A $g(q)=\frac{\Cov_{xq}}{\sqrt{\Var_q+\lambda}}\cdot\frac{q-\bar q}{\sqrt{\Var_q+\lambda}}+\bar x$: normalise $q$, then rescale and shift with statistics of $x$. $\lambda$ plays the role of $\epsilon$. Cost is about one LayerNorm; the normalisation structure supplies the extra gradient path.

\Q{Does $g$ add learnable parameters?}
\A No. $s_g$ and $b_g$ are computed in closed form from each block; only $\lambda$ is a fixed hyperparameter (0.01 everywhere).

\Q{State and prove Theorem~1.}
\A See Section~9. Centred rows/columns sum to zero, so cross terms vanish; the result is one integer matmul $Q_XQ_W$ plus rank-1 corrections.

\Q{The training code multiplies float reconstructions. Where is the saving?}
\A That is fake quantization, used so autograd can run. The reconstruction is affine with per-block scalars, so the scalars factor out of each dot product (Theorem~1). Deployed hardware does an integer matmul on $q_X,q_W$ and a cheap per-block rescale.

\Q{Why quantize activations, not just weights?}
\A Multiply cost scales with weight bits $\times$ activation bits. With float activations a 1-bit weight saves storage but the adds are still 16-bit. With both 1-bit the dot product is XNOR plus popcount. Activations also dominate memory traffic.

\Q{Do we compute $s_g,b_g$ for activations too? At what cost?}
\A Yes, per token, at every quantized layer (the paper uses affine $g$ for activations). This needs the float $x$ at run time: an $O(N)$ pass, about a LayerNorm. Weights' $s_g,b_g$ are precomputed. The energy proxy ignores this cost.

\Q{How is sparsity handled?}
\A Threshold small values to zero ($\delta_S$), quantize ($\delta_Q$), reconstruct with $g$ fitted to the dense $x$. $g$ adapts to the combined error and the backward pass sees both. For ternary sparse weights linear $g$ is used, which keeps zeros at zero.

\Q{What is 1.5-bit quantization?}
\A Three levels $\{-1,0,1\}$, which need $\log_23\approx1.58$ bits. At A1.5W1.5 STE gave NaN, $g$ reached 0.3297.

\Q{How do you get sub-1 bit per element?}
\A 1-bit values with 1:4 sparsity: one value plus a 2-bit position index per group of 4, so $(1+2)/4=0.75$ bits. This ignores the per-block scale and shift, which the paper does not count.

\Q{What is the energy proxy and how much should we trust it?}
\A Sparsity factor $\times$ activation bits $\times$ weight bits $\times$ ops. BF16 is $16\times16=256$, A1W1 is 1. It is an estimate: no hardware measurement, no memory-traffic or overhead model.

\Q{How does this compare with LSQ?}
\A LSQ (Esser et al., 2019) learns the step size by gradient descent, and its step-size gradient contains $\delta$. It does not change the upstream gradient to $x$. The new method computes scale and shift in closed form and its Jacobian with respect to $x$ contains $\delta$.

\Q{Are $s_f,b_f$ really constants?}
\A No, they depend on $x_{\max},x_{\min}$, and autograd flows through them. That only affects the two extreme entries per block. The paper's text and our equations hold them fixed; our Jacobian check excludes those entries.

\Q{Why did STE's activations blow up in your 1-bit run?}
\A At 1 bit STE reconstructs to $x_{\min}$ or $x_{\max}$, which overshoots most entries. Least squares puts the two levels near the group means. Over many layers STE's residual stream grew to about $340\times$ $g$'s.

\Q{Why didn't you reproduce the paper's loss explosion?}
\A Less compute: batch $32\times128$ and 1000--2000 steps versus $64\times256$ and 5000. We did reproduce the stall, the activation blow-up, and an upward loss turn at $3\times$ learning rate.

\Q{What is Tiny Shakespeare and what do the loss numbers mean?}
\A About 1.1M characters of Shakespeare dialogue; the model predicts the next character from a 65-character vocabulary. Loss in nats per character: random is $\ln65\approx4.17$; STE stalls near 2.6; $g$ reaches about 2.3; a float model reaches about 1.59.

\Q{Is the comparison in the paper fair?}
\A The STE vs.\ $g$ comparison is clean: same model, recipe and quantizer, only the reconstruction differs. Weaker points: single runs, no error bars, and BitNet run at A1W1 although its design uses 8-bit activations.

\Q{What is the main takeaway?}
\A A good forward approximation can still give a bad gradient. The rounding error looks ignorable, but training needs information about it. Fitting the reconstruction to the data puts that information into the backward pass.

\end{document}
